% !TEX program = xelatex
% 湖南科技大学本科生毕业设计（论文）LaTeX模板
% 纯模板版本 - 仅包含格式和结构

\documentclass[a4paper,12pt]{article}

% ==================== 基础包 ====================
\usepackage{geometry}
\usepackage{fontspec}
\usepackage{xeCJK}
\usepackage{setspace}
\usepackage{titlesec}
\usepackage{titletoc}
\usepackage{fancyhdr}
\usepackage{graphicx}
\usepackage{amsmath}
\usepackage{amssymb}
\usepackage{caption}
\usepackage{booktabs}
\usepackage{array}
\usepackage{longtable}
\usepackage[hidelinks]{hyperref}
\usepackage{cite}

% ==================== 页面设置 ====================
\geometry{
  a4paper,
  top=30mm,
  bottom=25mm,
  left=25mm,
  right=25mm,
  headheight=20mm,
  footskip=15mm
}

% ==================== 字体设置 ====================
\setCJKmainfont{SimSun}[AutoFakeBold=2.5,ItalicFont=KaiTi]
\setCJKsansfont{SimHei}[AutoFakeBold=2.5]
\setCJKmonofont{FangSong}
\setmainfont{Times New Roman}
\setsansfont{Arial}

% ==================== 页眉页脚设置 ====================
\pagestyle{fancy}
\fancyhf{}
\fancyhead[C]{\zihao{5}\songti 湖南科技大学本科生毕业设计（论文）}
\fancyfoot[C]{\zihao{5}-\thepage-}
\renewcommand{\headrulewidth}{0.4pt}

% ==================== 字号命令 ====================
\newcommand{\chuhao}{\fontsize{42pt}{\baselineskip}\selectfont}
\newcommand{\xiaochuhao}{\fontsize{36pt}{\baselineskip}\selectfont}
\newcommand{\yihao}{\fontsize{28pt}{\baselineskip}\selectfont}
\newcommand{\erhao}{\fontsize{21pt}{\baselineskip}\selectfont}
\newcommand{\xiaoerhao}{\fontsize{18pt}{\baselineskip}\selectfont}
\newcommand{\sanhao}{\fontsize{15.75pt}{\baselineskip}\selectfont}
\newcommand{\xiaosanhao}{\fontsize{15pt}{\baselineskip}\selectfont}
\newcommand{\sihao}{\fontsize{14pt}{\baselineskip}\selectfont}
\newcommand{\xiaosihao}{\fontsize{12pt}{\baselineskip}\selectfont}
\newcommand{\wuhao}{\fontsize{10.5pt}{\baselineskip}\selectfont}
\newcommand{\xiaowuhao}{\fontsize{9pt}{\baselineskip}\selectfont}

\newcommand{\zihao}[1]{%
  \ifnum#1=5\wuhao\fi
  \ifnum#1=-4\xiaosihao\fi
  \ifnum#1=4\sihao\fi
  \ifnum#1=-2\xiaoerhao\fi
  \ifnum#1=2\erhao\fi
  \ifnum#1=3\sanhao\fi
}

% ==================== 章节格式设置 ====================
\titleformat{\section}
  {\xiaoerhao\songti\bfseries\centering}
  {第\chinese{section}章}
  {1em}
  {}
\titlespacing*{\section}{0pt}{0.5\baselineskip}{0.5\baselineskip}

\titleformat{\subsection}
  {\sihao\songti\bfseries}
  {\arabic{section}.\arabic{subsection}}
  {1em}
  {}
\titlespacing*{\subsection}{0pt}{0.5\baselineskip}{0.5\baselineskip}

\titleformat{\subsubsection}
  {\xiaosihao\songti\bfseries}
  {\arabic{section}.\arabic{subsection}.\arabic{subsubsection}}
  {1em}
  {}
\titlespacing*{\subsubsection}{0pt}{0.5\baselineskip}{0.5\baselineskip}

% ==================== 正文格式设置 ====================
\setlength{\parindent}{2em}
\setstretch{1.25}
\xiaosihao\songti

% ==================== 图表标题格式 ====================
\captionsetup[figure]{
  labelsep=space,
  font={small,bf},
  name={图},
  labelfont=bf
}

\captionsetup[table]{
  labelsep=space,
  font={small,bf},
  name={表},
  labelfont=bf,
  position=top
}

% ==================== 目录格式设置 ====================
\titlecontents{section}[0em]
  {\xiaosihao\songti\bfseries}
  {第\chinese{section}章\quad}
  {}
  {\titlerule*[0.5pc]{.}\contentspage}

\titlecontents{subsection}[2em]
  {\xiaosihao\songti}
  {\thecontentslabel\quad}
  {}
  {\titlerule*[0.5pc]{.}\contentspage}

% ==================== 参考文献格式 ====================
\bibliographystyle{plain}
\renewcommand{\refname}{\sihao\songti\bfseries 参考文献}

% ==================== 摘要环境 ====================
\newenvironment{cnabstract}{
  \newpage
  \begin{center}
    \xiaoerhao\songti\bfseries 摘\quad 要
  \end{center}
  \vspace{1\baselineskip}
  \xiaosihao\songti
}{
  \vspace{1\baselineskip}
}

\newcommand{\cnkeywords}[1]{
  \noindent\textbf{关键词：}#1
}

\newenvironment{enabstract}{
  \newpage
  \begin{center}
    \xiaoerhao\bfseries ABSTRACT
  \end{center}
  \vspace{1\baselineskip}
  \xiaosihao
}{
  \vspace{1\baselineskip}
}

\newcommand{\enkeywords}[1]{
  \noindent\textbf{Keywords: }#1
}

% ==================== 文档信息（需填写）====================
\newcommand{\thesistitle}{论文题目}
\newcommand{\thesisauthor}{学生姓名}
\newcommand{\thesisstudentid}{学号}
\newcommand{\thesismajor}{专业}
\newcommand{\thesiscollege}{学院}
\newcommand{\thesisadvisor}{指导教师}
\newcommand{\thesisdate}{\today}

% ==================== 正文开始 ====================
\begin{document}

% ==================== 封面 ====================
\begin{titlepage}
  \centering
  \vspace*{2cm}

  {\erhao\heiti 湖南科技大学}

  \vspace{1cm}

  {\yihao\heiti 毕业设计（论文）}

  \vspace{3cm}

  \begin{tabular}{rl}
    {\sanhao\songti 题\quad 目：} & \underline{\makebox[8cm][c]{\sanhao\songti\thesistitle}} \\[1cm]
    {\sanhao\songti 学生姓名：} & \underline{\makebox[8cm][c]{\sanhao\songti\thesisauthor}} \\[1cm]
    {\sanhao\songti 学\quad 号：} & \underline{\makebox[8cm][c]{\sanhao\songti\thesisstudentid}} \\[1cm]
    {\sanhao\songti 专\quad 业：} & \underline{\makebox[8cm][c]{\sanhao\songti\thesismajor}} \\[1cm]
    {\sanhao\songti 学\quad 院：} & \underline{\makebox[8cm][c]{\sanhao\songti\thesiscollege}} \\[1cm]
    {\sanhao\songti 指导教师：} & \underline{\makebox[8cm][c]{\sanhao\songti\thesisadvisor}} \\[1cm]
  \end{tabular}

  \vfill

  {\sanhao\songti 二〇二\quad 年\quad 月\quad 日}
\end{titlepage}

% ==================== 扉页 ====================
\newpage
\thispagestyle{empty}
\begin{center}
  {\erhao\heiti 湖南科技大学}

  \vspace{0.5cm}

  {\xiaoerhao\heiti 毕业设计（论文）任务书}
\end{center}

\vspace{1cm}

\noindent
\begin{tabular}{ll}
  院 & 系（教研室）： \\[0.5cm]
  系（教研室）主任： & （签名）\quad\quad\quad\quad 年\quad 月\quad 日 \\[0.5cm]
\end{tabular}

\vspace{0.5cm}

\noindent
\begin{tabular}{lll}
  学生姓名： & 学号： & 专业： \\[0.5cm]
\end{tabular}

\vspace{0.5cm}

\noindent
1 设计（论文）题目及专题：

\vspace{2cm}

\noindent
2 学生设计（论文）时间：自\quad\quad 年\quad 月\quad 日开始至\quad\quad 年\quad 月\quad 日止

\vspace{1cm}

\noindent
3 设计（论文）所用资源和参考资料：

\vspace{3cm}

\noindent
4 设计（论文）应完成的主要内容：

\vspace{3cm}

\noindent
5 提交设计（论文）形式（设计说明与图纸或论文等）及要求：

\vspace{2cm}

\noindent
6 发题时间：\quad\quad 年\quad\quad 月\quad\quad 日

\vspace{1cm}

\noindent
\begin{tabular}{ll}
  指导教师： & （签名） \\[0.5cm]
  学\quad 生： & （签名） \\
\end{tabular}

% ==================== 指导人评语 ====================
\newpage
\thispagestyle{empty}
\begin{center}
  {\erhao\heiti 湖南科技大学}

  \vspace{0.5cm}

  {\xiaoerhao\heiti 毕业设计（论文）指导人评语}
\end{center}

\vspace{1cm}

\noindent
[主要对学生毕业设计（论文）的工作态度，研究内容与方法，工作量，文献应用，创新性，实用性，科学性，文本（图纸）规范程度，存在的不足等进行综合评价]

\vspace{8cm}

\noindent
\begin{tabular}{ll}
  指导人： & （签名） \\[0.5cm]
  & 年\quad 月\quad 日 \\[1cm]
  指导人评定成绩： & \\
\end{tabular}

% ==================== 评阅人评语 ====================
\newpage
\thispagestyle{empty}
\begin{center}
  {\erhao\heiti 湖南科技大学}

  \vspace{0.5cm}

  {\xiaoerhao\heiti 毕业设计（论文）评阅人评语}
\end{center}

\vspace{1cm}

\noindent
[主要对学生毕业设计（论文）的文本格式、图纸规范程度，工作量，研究内容与方法，实用性与科学性，结论和存在的不足等进行综合评价]

\vspace{8cm}

\noindent
\begin{tabular}{ll}
  评阅人： & （签名） \\[0.5cm]
  & 年\quad 月\quad 日 \\[1cm]
  评阅人评定成绩： & \\
\end{tabular}

% ==================== 答辩记录 ====================
\newpage
\thispagestyle{empty}
\begin{center}
  {\erhao\heiti 湖南科技大学}

  \vspace{0.5cm}

  {\xiaoerhao\heiti 毕业设计（论文）答辩记录}
\end{center}

\vspace{1cm}

\noindent
\begin{tabular}{ll}
  日期： & \\[0.5cm]
  学生： & 学号：\quad\quad\quad\quad 班级： \\[0.5cm]
  题目： & \\[0.5cm]
\end{tabular}

\vspace{0.5cm}

\noindent
提交毕业设计（论文）答辩委员会下列材料：

\vspace{0.5cm}

\noindent
\begin{tabular}{ll}
  1 设计说明书/ 论 文 & 共\quad\quad 页 \\[0.3cm]
  2 设计（论文）图 纸 & 共\quad\quad 页 \\[0.3cm]
  3 指导人、评阅人评语 & 共\quad\quad 页 \\
\end{tabular}

\vspace{1cm}

\noindent
毕业设计（论文）答辩委员会评语：

\vspace{0.5cm}

\noindent
[主要对学生毕业设计（论文）的研究思路，设计（论文）质量，文本图纸规范程度和对设计（论文）的介绍，回答问题情况等进行综合评价]

\vspace{5cm}

\noindent
\begin{tabular}{ll}
  答辩委员会主任： & （签名） \\[0.3cm]
  委员： & （签名） \\[0.3cm]
  & （签名） \\[0.3cm]
  & （签名） \\[0.3cm]
  & （签名） \\[0.5cm]
  答辩成绩： & \\[0.5cm]
  总评成绩： & \\
\end{tabular}

% 前置部分使用罗马数字页码
\newpage
\pagenumbering{roman}

% ==================== 中文摘要 ====================
\begin{cnabstract}
（此处填写中文摘要内容，小四号宋体，1.25倍行距）

\cnkeywords{关键词1；关键词2；关键词3；关键词4；关键词5}
\end{cnabstract}

% ==================== 英文摘要 ====================
\begin{enabstract}
(Abstract content in English, font size: 12pt, line spacing: 1.25)

\enkeywords{Keyword1; Keyword2; Keyword3; Keyword4; Keyword5}
\end{enabstract}

% ==================== 目录 ====================
\newpage
\tableofcontents

% 正文部分使用阿拉伯数字页码
\newpage
\pagenumbering{arabic}

% ==================== 第一章 ====================
\section{章标题}

\subsection{节标题}

正文内容（小四号宋体，1.25倍行距，首行缩进2字符）

\subsubsection{条标题}

正文内容

% ==================== 第二章 ====================
\section{章标题}

\subsection{节标题}

正文内容

% ==================== 参考文献 ====================
\newpage
\begin{thebibliography}{99}
\addcontentsline{toc}{section}{参考文献}

\bibitem{ref1}
作者. 文献题名[文献类型标识]. 出版地: 出版者, 出版年: 起止页码.

\end{thebibliography}

% ==================== 致谢 ====================
\newpage
\section*{致\quad 谢}
\addcontentsline{toc}{section}{致谢}

（此处填写致谢内容）

% ==================== 附录 ====================
\newpage
\section*{附录A：附录标题}
\addcontentsline{toc}{section}{附录A：附录标题}

（此处填写附录内容）

\end{document}
